% Generated draft: requires mathematical and source review.
\documentclass[11pt]{article}
\usepackage[margin=25mm]{geometry}
\usepackage{fontspec}
\setmainfont{Latin Modern Roman}
\usepackage{amsmath,amssymb,mathtools,unicode-math}
\setmathfont{Latin Modern Math}
\usepackage{enumitem,needspace,xurl,hyperref}
\setlength{\parindent}{0pt}
\setlength{\parskip}{5pt}
\setlength{\emergencystretch}{3em}
\newcommand{\N}{\mathbb{N}}
\newcommand{\Z}{\mathbb{Z}}
\newcommand{\Q}{\mathbb{Q}}
\newcommand{\R}{\mathbb{R}}
\newcommand{\C}{\mathbb{C}}
\newcommand{\F}{\mathbb{F}}
\newcommand{\A}{\mathbb{A}}
\newcommand{\PP}{\mathbb{P}}
\newcommand{\E}{\mathbb{E}}
\newcommand{\eps}{\varepsilon}
\newcommand{\dd}{\,\mathrm{d}}
\begin{document}
\section*{UnsolvedMath ID: 30001693}
\textbf{Source problem code:} catalogue entry
\subsection{Definitions and mathematical statement}
Fix integers \(d\geq 1\) and \(r\geq 2\). For a positive integer \(n\), let \(\Delta^{n-1}\) be the simplex with vertex set \(V\), where \(|V|=n\). A coloring is a partition
\[
V=C_1\sqcup\cdots\sqcup C_{d+1}.
\]
A nonempty face \(\sigma\) is \emph{rainbow} if its vertex set contains at most one vertex from each \(C_i\). Faces are pairwise disjoint when their vertex sets are pairwise disjoint. The required intersection property is the existence of nonempty, pairwise disjoint rainbow faces \(\sigma_1,\ldots,\sigma_r\) such that
\[
\bigcap_{j=1}^{r}f(\sigma_j)\neq\varnothing.
\]
Here \(f:\Delta^{n-1}\to\R^d\) is affine; no general-position or injectivity assumption is supplied.

For \(n\geq(d+1)r\), let \(P_{d,r}(n)\) mean that for every partition of \(V\) into \(d+1\) color classes satisfying \(|C_i|\geq r\) for every \(i\), and every affine map \(f:\Delta^{n-1}\to\R^d\), the required intersection property holds. Define the feasible version of the supplied minimum by
\[
n_{\mathrm{feas}}(d,r)=\min\{n\in\Z:n\geq(d+1)r\text{ and }P_{d,r}(n)\},
\]
provided this set is nonempty. The first requested determination is this minimum, with conjectured value
\[
n_{\mathrm{feas}}(d,r)=(d+1)r.
\]
\emph{Unresolved qualification:} the source writes \(n(d,r)\) without explicitly restricting the minimum to feasible values of \(n\). For \(n<(d+1)r\), no prescribed partition exists, so the literal universal assertion is vacuous. The notation \(n_{\mathrm{feas}}\) makes this qualification explicit rather than silently identifying it with the source's undefined minimum.

For a positive integer \(t\), let \(Q_{d,r}(t)\) mean that for every partition of the vertices of \(\Delta^{(d+1)t-1}\) into exactly \(d+1\) color classes of size \(t\), and every affine map
\[
f:\Delta^{(d+1)t-1}\to\R^d,
\]
the required intersection property holds. Define
\[
t(d,r)=\min\{t\in\Z:t\geq1\text{ and }Q_{d,r}(t)\},
\]
provided this set is nonempty. The alternative requested determination is \(t(d,r)\), with conjectured value
\[
t(d,r)=r.
\]
The supplied record calls this equivalent to the B\'ar\'any--Larman equality. At the proposed minimum \(n=(d+1)r\), the condition \(|C_i|\geq r\) forces every color class to have size exactly \(r\), so the two proposed assertions have the same intersection requirement there. The feasibility qualification for the first minimum remains necessary.
\subsection{Short English statement}
For every \(d\geq1\) and \(r\geq2\), determine the smallest feasible number of colored vertices that guarantees \(r\) vertex-disjoint rainbow faces with intersecting affine images in \(\R^d\). The proposed value is \((d+1)r\). With exactly \(d+1\) equally sized color classes, determine the smallest class size; the proposed value is \(r\). The source's first minimum requires an explicit feasibility restriction to avoid vacuous cases.
\subsection{Sources}
\begin{itemize}
\item Local UnsolvedMath registry, permanent ID \(30001693\), title ``B\'ar\'any--Larman Colored Tverberg Conjecture,'' source problem code \texttt{catalogue entry}, in the Oberwolfach Reports Open Problems collection. The record gives proposed year \(2011\) and does not name an individual proposer.
\item Supplied report citation: \emph{Topological and Geometric Combinatorics} (2011), Oberwolfach Reports, pp.\ 349--423, DOI \texttt{10.4171/owr/2011/08}. Supplied literature-source label: ``original OWR report.'' \url{https://doi.org/10.4171/owr/2011/08}. The background describes these reports as documenting mathematical workshops and their open-problem discussions.
\item Supplied literature-source label: ``optimal colored Tverberg bounds.'' \url{https://arxiv.org/abs/0910.4987}. No authors or full bibliographic title are supplied in the record.
\item Literature assessment dated \texttt{2026-08-21}: supplied dated curation, not independently verified. The record labels the problem \texttt{partially\_solved} and states that the B\'ar\'any--Larman equality is proved in low dimensions and when \(r+1\) is prime, that asymptotically optimal general bounds are available, and that the full equality remains open. The dimensions and bounds are not specified. The background expressly asks for independent verification of time-sensitive status and cited claims.
\item The registry labels the statement \texttt{corrected\_verified} and reports a manual semantic, typographic, and LaTeX review of its cleanup. These are supplied provenance claims, not independent mathematical verification of this draft or its literature status.
\end{itemize}
\end{document}
